% Part: proof-theory
% Chapter: cut-elimination
% Section: introduction

\documentclass[../../../include/open-logic-section]{subfiles}

\begin{document}

\olfileid{pt}{cut}{int}

\olsection{مقدمه}

قاعدهٔ \Cut{} قاعدهٔ زیر است:
\[
\Axiom$\Gamma \fCenter \Delta, !A$
\Axiom$!A, \Pi \fCenter \Lambda$
\RightLabel{\Cut}
\BinaryInf$\Gamma, \Pi \fCenter \Delta, \Lambda$
\DisplayProof
\]
(!!{formula}~$!A$ را \emph{فرمول برش} می‌نامند.)

\Cut{} در حساب دنباله‌ای اصلی گنتسن،~\Log{LK}، وجود داشت. \emph{قضیهٔ
اصلی} گنتسن (\emph{Hauptsatz}) می‌گوید هر دنبالهٔ !!{provable} در
~\Log{LK} بدون استفاده از قاعدهٔ \Cut{} نیز !!{provable} است؛ یعنی
می‌توان \Cut{} را از \Log{LK}-برهان‌ها «حذف» کرد. دستگاه‌های
\Log{G1c} و \Log{G3c} ما \Cut را در بر ندارند، اما همین پرسش را
می‌توان دربارهٔ آن‌ها نیز مطرح کرد: آیا می‌توان \Cut{} را از
!!{proof}هایی حذف کرد که \Cut{} را همراه با دیگر قواعد \Log{G1c}
(یا \Log{G3c}) به کار می‌برند؟ در اصطلاحی که پیش‌تر تثبیت کردیم،
پرسش هم‌ارز این است: آیا \Cut{} قاعده‌ای پذیرفتنی در \Log{G1c}
(یا \Log{G3c}) است؟

قضیهٔ حذف برش شاید مهم‌ترین و مشهورترین نتیجهٔ نظریهٔ اثبات باشد.
نخست ثابت می‌کند \Log{G1c} و \Log{G3c} به قاعده‌ای نیاز ندارند که
ممکن است در نگاه اول آشکارا ضروری به نظر برسد. وقتی برهان‌های واقعی
را صوری می‌کنید، درمی‌یابید باید راهی برای استفاده از لم‌ها داشته
باشید. ریاضی‌دانان دوست دارند نتیجه‌هایی را اثبات کنند و سپس در
برهان نتایج دیگر به آن‌ها استناد کنند. ممکن است نخست برهان نتیجهٔ
$L$، یعنی لم، را به صورت !!{proof} دنبالهٔ $\Sequent L$ صوری کنید.
سپس برهان نتیجهٔ دوم~$R$ را که از~$L$ استفاده می‌کند، به صورت
!!{proof} دنبالهٔ $L \Sequent R$ صوری می‌کنید. از کجا می‌دانید
برهانی برای~$R$، یعنی !!a{proof} برای صرفاً $\Sequent R$، وجود دارد؟
باید بتوانید دو !!{proof} را ترکیب کنید، و قاعدهٔ \Cut{} چنین امکانی
می‌دهد.

گفتیم \Log{G1c} و \Log{G3c} تمام‌اند و بنابراین هر دنباله‌ای را که
انتظار می‌رود اثبات‌پذیر باشد، یعنی هر دنبالهٔ معتبری را، اثبات
می‌کنند. این امر را می‌توان مستقیم اثبات کرد. اما زمانی که گنتسن
می‌نوشت، تنها تمامیت دستگاه‌های استنتاج اصل‌موضوعی شناخته شده بود.
گنتسن برای نشان‌دادن تمامیت \Log{LK} ترجمه‌ای از !!{derivation}های
دستگاه اصل‌موضوعی به \Log{LK}-!!{proof}ها ارائه کرد. این ترجمه به‌طور
اساسی از \Cut{} استفاده می‌کرد. حذف برش فقط نتیجه‌ای مهم برای منطق
کلاسیک نیست: بسیاری منطق‌های دیگر نیز حساب دنباله‌ای دارند. برای برخی
از آن‌ها، اثبات مستقیم تمامیت حساب دنباله‌ای دشوار یا ناممکن است؛ پس
ترجمه از دستگاه‌های دیگر، با استفاده از \Cut، تنها راه نشان‌دادن
تمامیت است. آنگاه اثبات نظریه‌اثباتی حذف برش لازم است تا ثابت شود
\Cut{} زائد است و دستگاه بدون \Cut{} تمام است.

حذف برش از آن رو نیز مهم است که می‌توان آن را برای به‌دست‌آوردن
نتایج مستقل و جالب دیگر به کار برد. دو موردی که بررسی خواهیم کرد لم
درونیابی کریگ و قضیهٔ هربراند هستند.

ایده‌های متداول در اثبات حذف برش نشان می‌دهند چگونه !!a{proof}ی که
از \Cut{} استفاده می‌کند به برهانی بدون آن تبدیل می‌شود. این تبدیل‌ها
معمولاً «موضعی»اند؛ یعنی بخشی از !!a{proof}، معمولاً یک
زیر-!!{proof} که به استنتاج \Cut{} پایان می‌یابد، را برمی‌داریم و با
زیر-!!{proof} دیگری از همان دنباله جایگزین می‌کنیم که در آن استنتاج
\Cut{} حذف شده یا با استنتاج‌های دیگری جایگزین شده است. چنین
استدلال‌هایی الگوریتم‌های گام‌به‌گامی برای تبدیل !!{proof}های دارای
\Cut{} به برهان‌های بدون آن فراهم می‌کنند. این ویژگی هنگام کاربرد
روش‌های حذف برش اطلاعات بیشتری می‌دهد. برای نمونه، برهان‌های
نظریه‌مدلی‌ای برای لم درونیابی هست که می‌گویند «اگر
$!A \lif !B$ معتبر باشد، آنگاه یک درونیاب وجود دارد» (فعلاً مهم نیست
درونیاب چیست). یک استدلال نظریه‌اثباتی با استفاده از حذف برش الگوریتمی
فراهم می‌کند که !!a{proof}ی برای $!A \lif !B$ می‌گیرد و یک درونیاب
بازمی‌گرداند. برخلاف استدلال نظریه‌مدلی، استدلال نظریه‌اثباتی
\emph{سازنده} است.

دو راه رایج برای اثبات قضیهٔ حذف برش وجود دارد. راه نخست، که به گنتسن
بازمی‌گردد، نشان می‌دهد می‌توان بالاترین برش‌ها را از یک برهان حذف کرد
و سپس آن‌ها را یکی‌یکی حذف می‌کند تا هیچ‌یک باقی نماند. راه دوم، که
نخست شوتّه و تیت ارائه کردند، بر پیچیده‌ترین برش‌های یک !!a{proof}
تمرکز می‌کند و پی‌درپی برش‌هایی را حذف می‌کند که بیشینه‌اند؛ یعنی هیچ
برش دیگری !!{formula} برشی با عمق بیشتر ندارد. هر روش مزایا و معایبی دارد.
برای برخی دستگاه‌ها یک روش کار می‌کند، حال آنکه اجرای روش دیگر دشوار
یا حتی ناممکن است. بنابراین هر دو را شرح می‌دهیم. حذف برش را برای
~\Log{G3c} اثبات خواهیم کرد. در واقع نه حذف \Cut{}، بلکه حذف
\emph{برش با بافت مشترک}~\CutCS را نشان خواهیم داد:
\[
\Axiom$\Gamma \fCenter \Delta, !A$
\Axiom$!A, \Gamma \fCenter \Delta$
\RightLabel{\CutCS}
\BinaryInf$\Gamma \fCenter \Delta$
\DisplayProof
\]
اگر یک حساب دنباله‌ای تضعیف و انقباض را مجاز بداند، خواه به صورت
قواعد واقعی مانند \Log{G1c} و خواه به صورت قواعد پذیرفتنی مانند
\Log{G3c}، \Cut{} می‌تواند \CutCS را شبیه‌سازی کند و برعکس. \CutCS{}
را به جای \Cut برمی‌گزینیم، زیرا راهبرد نخست حذف برش برای~\CutCS
آسان‌تر بیان می‌شود و راهبرد دوم فقط برای~\CutCS کار می‌کند.

\end{document}
